\documentclass[11pt]{amsart}
\usepackage[T1]{fontenc}
\usepackage{lmodern,amsmath,amssymb,amsthm,mathtools,mathrsfs,microtype}
\usepackage[margin=1in]{geometry}
\usepackage[hidelinks]{hyperref}
\numberwithin{equation}{section}
\newtheorem{theorem}{Theorem}
\newtheorem{conjecture}{Conjecture}
\newtheorem{lemma}{Lemma}[section]
\newtheorem{proposition}[lemma]{Proposition}
\newcommand{\HH}{\mathscr H}
\newcommand{\T}{\mathbb T}
\newcommand{\SN}{S_{\mathrm N}}
\newcommand{\OO}{\mathcal O}
\newcommand{\TT}{\mathcal T}
\newcommand{\KK}{\mathcal K}
\newcommand{\NN}{\mathcal N}
\newcommand{\Aop}{\mathcal A}
\newcommand{\PP}{\mathcal P}
\newcommand{\JJ}{\mathcal J}
\newcommand{\VV}{\mathcal V}
\newcommand{\tr}{\operatorname{tr}}
\newcommand{\ran}{\operatorname{ran}}
\newcommand{\Dom}{\operatorname{Dom}}
\newcommand{\indm}{\operatorname{ind}_{-}}
\newcommand{\ip}[2]{\langle #1,#2\rangle}
\newcommand{\norm}[1]{\lVert #1\rVert}
\newcommand{\Argm}{\operatorname{Arg}_{-}}
\newcommand{\ph}{\varphi}
\title[The Neumann P\'olya inequality]{The Neumann P\'olya inequality on simply connected planar domains}
\date{}
\hypersetup{pdftitle={The Neumann Polya inequality on simply connected planar domains},pdfsubject={Neumann eigenvalues on simply connected planar domains}}
\begin{document}
\maketitle

\section{Introduction}\label{sec:intro}

Let $\Omega\subset\mathbb R^2$ be a bounded connected Lipschitz domain,
and let $|\Omega|$ denote its area. The Neumann Laplacian
$A_N=-\Delta_\Omega^N$ is the self-adjoint operator associated with the
quadratic form
\[
 q[u]=\int_\Omega|\nabla u|^2,\qquad u\in H^1(\Omega).
\]
Its eigenvalues, repeated according to multiplicity, are denoted by
\[
 0=\mu_0(\Omega)<\mu_1(\Omega)\le\mu_2(\Omega)\le\cdots.
\]
We include the zero eigenvalue in the strict counting function
\[
 N_N(E)=\#\{j\ge0:\mu_j(\Omega)<E\}.
\]
Weyl's law identifies the leading term as
$N_N(E)\sim |\Omega|E/(4\pi)$ as $E\to\infty$; see
\cite{Weyl,FLPS}. P\'olya's conjecture asks whether this asymptotic
expression is a lower bound at every positive energy.

The conjecture goes back to P\'olya's discussion of vibrating membranes
in \cite{Polya54}. His subsequent paper \cite{Polya} established the
corresponding inequalities for tiling domains, with an additional
restriction in the Neumann case that was removed by Kellner
\cite{Kellner}. In the indexing used here, the planar Neumann part of
the conjecture takes the following form.

\begin{conjecture}\label{conj:polya}
For every bounded connected Lipschitz domain
$\Omega\subset\mathbb R^2$ and every $j\ge1$,
\[
 |\Omega|\mu_j(\Omega)\le4\pi j.
\]
\end{conjecture}

The same problem can be stated as $N_N(E)\ge |\Omega|E/(4\pi)$ for
all $E>0$. The strict version, in which the eigenvalue inequality is
replaced by $|\Omega|\mu_j(\Omega)<4\pi j$, is called the strong
P\'olya conjecture in \cite{FLP}. The question concerns the whole
spectrum; an asymptotic formula alone does not settle the low- and
intermediate-energy ranges.

Sharp geometric estimates for the first positive Neumann eigenvalue
were obtained by Szeg\H{o} \cite{Szego} for simply connected planar
domains and by Weinberger \cite{Weinberger} without that topological
restriction and in arbitrary dimension. For the full spectrum,
Kr\"oger \cite{Kroger} proved sharp bounds for eigenvalue sums and
upper bounds for individual eigenvalues with a larger constant.
Laptev \cite{Laptev} developed the associated trace inequalities and
product-domain arguments. These results do not give the conjectured
constant for individual eigenvalues on an arbitrary planar domain.

Beyond tiling domains, Filonov, Levitin, Polterovich, and Sher
\cite{FLPS} proved the conjecture for the disk and for arbitrary planar
sectors, as well as its Dirichlet version for balls in all dimensions.
Freitas and Salavessa \cite{FS} obtained further families of non-tiling
domains by using thin-domain constructions. He and Wang \cite{HW}
proved the conjecture for sufficiently thin product domains.
More recently, Filonov
\cite{Filonov} improved the universal Neumann counting bound for planar
convex domains. Frank and Larson \cite{FrankLarson} extended sharp
semiclassical inequalities on convex domains to a range of Riesz
exponents below one. Here a Riesz mean of positive order $\alpha$ is
$\sum_{j\ge0}(E-\mu_j)_+^\alpha$, where $t_+=\max\{t,0\}$;
bounds for such means are distinct from the order-zero counting
inequality. The higher-dimensional Neumann ball case was established
in \cite{FLPS26}; see also \cite{Li26}. None of these results covers all simply connected
planar Lipschitz domains.

We prove the strong Neumann conjecture for this class.

\begin{theorem}\label{thm:main}
For every bounded simply connected Lipschitz domain
$\Omega\subset\mathbb R^2$ and every $j\ge1$,
\begin{equation}\label{eq:main}
 |\Omega|\mu_j(\Omega)<4\pi j.
\end{equation}
\end{theorem}

Theorem~\ref{thm:main} establishes Conjecture~\ref{conj:polya}, with
strict inequality, under the additional assumption of simple
connectedness. Equivalently,
\[
 N_N(E)>\frac{E|\Omega|}{4\pi}\qquad(E>0).
\]
Indeed, the counting inequality at $E=\mu_j$ gives
\eqref{eq:main}, because $N_N(\mu_j)\le j$. Conversely, if
$m=N_N(E)$, then $m\ge1$ and $E\le\mu_m$, so \eqref{eq:main}
gives $E|\Omega|<4\pi m$.

\subsection{Formal verification}

Theorem~\ref{thm:main} and its equivalent strict counting formulation have
been formalized and machine-checked in Lean 4.28.0.  The development is
available at
\url{https://github.com/QuanyuTang/polya-neumann-simply-connected-formalization}.
It identifies $\mathbb R^2$ with $\mathbb C$, defines the Neumann
eigenvalues by the variational min--max construction, and proves the stated
inequalities for bounded simply connected Lipschitz domains.  The repository
keeps the proof modular and provides a pinned Lean/Mathlib build together
with an axiom audit for the main theorem and its principal intermediate
results.  The audited declarations use only Lean's standard
$\operatorname{propext}$, $\operatorname{Classical.choice}$, and
$\operatorname{Quot.sound}$ axioms; no \texttt{sorry} or
\texttt{native\_decide} proof step
is used.

The proof uses a unitary evolution along the boundary. Its monodromy,
defined in \eqref{eq:transport}, is denoted by $V_E$, and we set
$U_E=V_E^*$. For a smooth Jordan domain, $U_E-I$ is trace class and
its Fredholm determinant is exactly $\exp(iE|\Omega|/2)$. The
Neumann-to-Dirichlet map relates a Cayley transform of $V_E$ to a
boundary quadratic form. In conformal coordinates, the principal
part of that form vanishes on the nonpositive Fourier modes. A
transmutation of antiholomorphic functions, based on the method of
Vekua \cite{Vekua,MHP}, identifies the corresponding kernel exactly.
After a finite-rank correction at the cut in the boundary coordinate,
this gives a uniform one-sided estimate for the Cayley transform.
At a Neumann eigenvalue of multiplicity $m$, at most $m$ dimensions
are excluded from that estimate.

To describe the counting step, call an energy \emph{nonresonant} if
it is not a Neumann eigenvalue. At such an energy let $\Phi_-(E)$
be the sum, with multiplicities, of the arguments in $(-2\pi,0)$ of
the eigenvalues of $U_E$ other than $1$. Absolute convergence is proved
in Section~\ref{sec:counting}. The determinant identity makes
\begin{equation}\label{eq:integer-intro}
 M(E)=\frac{E|\Omega|/2-\Phi_-(E)}{2\pi}
\end{equation}
an integer. The one-sided estimate shows that $M$ is constant between
Neumann eigenvalues and that its increase across an eigenvalue is at
most its multiplicity. A separate small-energy calculation gives
$M=1$ below the first positive Neumann eigenvalue. The comparison
$M(E)\le N_N(E)$ then yields
\[
 N_N(E)\ge\frac{E|\Omega|}{4\pi}
             +\frac{\|U_E-I\|}{2\pi}.
\]
The extra term is positive. Its continuity under boundary
approximation is what preserves strictness for Lipschitz domains.
The use of eigenphases has a precedent in inside--outside duality
\cite{EP}, but the boundary evolution here is not identified with an
exterior scattering matrix. In particular, the argument does not
require continuation of individual eigenvectors through the
essential spectral point $1$.

Section~\ref{sec:transport} constructs the evolution and establishes
its relation to Helmholtz Cauchy data. Section~\ref{sec:transmutation}
proves the transmutation and kernel results.
Section~\ref{sec:periodic} treats the periodic boundary form, and
Section~\ref{sec:uniform} proves the uniform estimate at nonresonant
and resonant energies. Sections~\ref{sec:small} and
\ref{sec:counting} give the normalization and the counting argument.
Finally, Section~\ref{sec:lipschitz} passes to Lipschitz boundaries
using the smooth approximations of Ball and Zarnescu \cite{BZ}.

\section{Boundary transport and Cauchy data}\label{sec:transport}

All Hilbert spaces are complex, and inner products are conjugate-linear
in the first variable. We identify $\mathbb R^2$ with $\mathbb C$.
For a bounded operator, an unqualified norm is its operator norm.
Boundary coordinate measure is $d\theta$ on $\T=\mathbb R/(2\pi\mathbb Z)$;
direction measure is $d\phi/(2\pi)$. Sobolev duality extends the $L^2$
pairing. On boundary Fourier modes $\ph_n=(2\pi)^{-1/2}e^{in\theta}$ put
\[
 D\ph_n=n\ph_n,\quad \Lambda\ph_n=(1+|n|)\ph_n,
 \quad \mathcal H_+=\overline{\operatorname{span}}\{\ph_n:n>0\},
 \quad \mathcal H_{\le0}=\mathcal H_+^\perp.
\]
Set $\mathcal H_-=\overline{\operatorname{span}}\{\ph_n:n<0\}$.
The orthogonal projections onto $\mathcal H_+$, $\mathcal H_-$, and
$\mathbb C\ph_0$ are $\Pi_+$, $\Pi_-$, and $\Pi_0$, respectively;
$\Pi_{\le0}=\Pi_-+\Pi_0$ and $D_+=D|_{\mathcal H_+}$.
We use $\norm{g}_{H^s}=\norm{\Lambda^s g}_{L^2}$ and set
$X=H^{-1/2}(\T)$, $X^*=H^{1/2}(\T)$. The symbols $\mathcal S_1$
and $\mathcal S_2$ denote the trace-class and Hilbert--Schmidt ideals;
$\det_F$ denotes the Fredholm determinant of an identity plus a
trace-class operator.

We set $D^{-1}\ph_0=|D|^{-1}\ph_0=0$ when these inverses occur.
The spectrum and point spectrum of an operator are denoted by $\sigma$
and $\sigma_p$, respectively. For a self-adjoint operator $H$,
$\mathbf1_I(H)$ is its spectral projection on a Borel set $I$.
The negative index $\indm H$ is the supremum of the dimensions of
subspaces on which its quadratic form is strictly negative.
The boundary trace is denoted by $\operatorname{Tr}$.
We write $\tr$ for the operator trace. All estimates that are locally
uniform in energy are for the fixed smooth domain under consideration.

For now $\Omega$ is a bounded $C^\infty$ Jordan domain. Choose a positive
smooth regular boundary parametrization $\gamma$ on $[0,L]$, where $L=2\pi$ is the
parameter period, not the boundary length,
and translate so that $\gamma(0)=\gamma(L)=0$. On
$\HH=\ell^2(\mathbb N_0)$, with standard orthonormal basis
$(e_n)_{n\ge0}$, let
\[
 \SN e_0=\sqrt2 e_1,\qquad \SN e_n=e_{n+1}\ (n\ge1),
 \qquad [\SN^*,\SN]=2P_0-P_1=:D_N.
\]
Here $\mathbb N_0=\{0,1,2,\ldots\}$ and $P_n$ is projection onto $\mathbb C e_n$. For $E=k^2>0$ solve
\begin{equation}\label{eq:transport}
 W_E'=-\frac{ik}{2}(\gamma'\SN^*+\overline{\gamma'}\SN)W_E,
 \quad W_E(0)=I,\qquad V_E=W_E(L),\quad U_E=V_E^*.
\end{equation}
The coefficient in \eqref{eq:transport} is bounded and skew-adjoint,
so $W_E(\theta)$ is unitary. At $E=0$ set
$W_0(\theta)=V_0=U_0=I$.
We abbreviate the conormal density by
$\partial_\nu^\gamma u=|\gamma'|(\partial_\nu u)\circ\gamma$,
where $\nu$ is the outward unit normal. Define
\begin{align*}
 (\OO_Ev)(\theta)&=\ip{e_0}{W_E(\theta)v},\\
 (\TT_Eg)(\theta)&=\int_0^\theta
 \ip{e_0}{W_E(\theta)W_E(s)^*e_0}g(s)\,ds.
\end{align*}
Initially $\OO_E:\HH\to L^2(0,L)$ and
$\TT_E:L^2(0,L)\to H^1(0,L)$. Their values are functions on the
cut interval, not necessarily periodic. Directly,
\begin{equation}\label{eq:OTadj}
 \OO_E^*g=\int_0^L W_E(s)^*e_0g(s)\,ds,
 \qquad \TT_E+\TT_E^*=\OO_E\OO_E^*,
 \qquad\norm{\OO_E}_{\mathcal S_2}^2=L.
\end{equation}

\begin{lemma}\label{lem:area}
The path $U_E-I$ is trace-norm continuously differentiable for $E\ge0$
and trace-norm analytic for $E>0$. If $L_E=-iU_E^*\partial_EU_E$, then
\begin{equation}\label{eq:det}
 \tr L_E=|\Omega|/2,\qquad
 \det\nolimits_F U_E=\exp(iE|\Omega|/2),\qquad
 L_0=\tfrac12|\Omega|D_N.
\end{equation}
\end{lemma}
\begin{proof}
Set $G_x=-i(\SN+\SN^*)/2$ and $G_y=(\SN^*-\SN)/2$; then
$[G_x,G_y]=iD_N/2$. Writing $\gamma=x+iy$, use the two coefficients
$C_\theta=\sqrt E(x_\theta G_x+y_\theta G_y)$ and
$\mathcal C_E=(xG_x+yG_y)/(2\sqrt E)$. Differentiation of the evolution gives
\[
 \partial_\theta(W_E^*\partial_EW_E-W_E^*\mathcal C_EW_E)
 =W_E^*(\partial_EC_\theta-\partial_\theta \mathcal C_E+[C_\theta,\mathcal C_E])W_E.
\]
The mixed derivatives cancel, the endpoint values of $\mathcal C_E$ vanish, and
$[C_\theta,\mathcal C_E]=-ibD_N/4$, where
$b=\operatorname{Im}(\overline\gamma\gamma')$. Hence
\begin{equation}\label{eq:curvature}
 V_E^*\partial_EV_E=-iQ_E,\qquad
 Q_E=\frac14\int_0^L b(\theta)W_E(\theta)^*D_NW_E(\theta)\,d\theta,
 \qquad L_E=V_EQ_EV_E^*.
\end{equation}
The integral is a trace-norm integral of finite-rank operators and
$\norm{Q_E}_{\mathcal S_1}\le3\int|b|/4$. Since
$\int b=2|\Omega|$ and $\tr D_N=1$, it has the asserted trace and limit
at zero. The Picard series gives analyticity for positive energies and
continuity at zero. Integrate $\partial_EV_E=-iV_EQ_E$ in trace norm
from zero. The logarithmic derivative of the Fredholm determinant is
$i\tr L_E$, which proves \eqref{eq:det}; see also the standard determinant
calculus in \cite{Simon}. No sign of $b$ has been assumed.
\end{proof}

For direction density $a\in L^2(\T,d\phi/(2\pi))$ write
\[
 u_a(z)=\int_0^{2\pi}a(\phi)e^{-ik\operatorname{Re}(ze^{-i\phi})}
                   \frac{d\phi}{2\pi},\qquad
 h_a=u_a\circ\gamma,\quad g_a=|\gamma'|\partial_\nu u_a\circ\gamma.
\]
Put $\partial=(\partial_x-i\partial_y)/2$ and
$\bar\partial=(\partial_x+i\partial_y)/2$. With
$y_0=u_a/\sqrt2$ and $y_n=(2i/k)^n\partial^n u_a$ for $n\ge1$,
Parseval gives a vector in $\HH$ with legitimate Hilbert-space derivatives.
Since the outward normal is $-i\gamma'/|\gamma'|$,
\begin{equation}\label{eq:driven}
 y'=C_\theta y-\frac i{\sqrt2}g_ae_0,
 \qquad h_a'= -ig_a-ik\gamma'y_1.
\end{equation}
If $v=y(0)$, variation of constants and periodicity give
\begin{align}
 h_a&=\sqrt2\OO_Ev-i\TT_Eg_a,\label{eq:trace-driven}\\
 (I-V_E)v&=-\frac i{\sqrt2}V_E\OO_E^*g_a,
 \qquad \OO_E^*g_a=\sqrt2 i(V_E^*-I)v.\label{eq:endpoint}
\end{align}
Every $v\in\HH$ occurs: if
$a(\phi)=\sum_{n\in\mathbb Z}a_ne^{in\phi}$, take
$a_0=\sqrt2v_0$, $a_n=v_n$ for $n\ge1$, and $a_n=0$ for $n<0$.
These coefficients belong to $\ell^2(\mathbb Z)$, so they define an
$L^2$ density.

We use the layer-potential mapping and jump relations, Rellich's
exterior uniqueness theorem, and uniqueness for zero interior Cauchy
data; see \cite{McLean}. Approximation by Herglotz waves is treated
in \cite{CK}. The boundary-density statement needed here includes
resonant energies and the $H^{-1/2}$ topology, so we give its
layer-potential proof. We use the convention in which the interior
trace minus the exterior trace of a double-layer potential with
density $h$ is $-h$, whereas the corresponding jump of the normal
derivative of a single-layer potential with density $g$ is $g$.

\begin{lemma}\label{lem:Cauchy}
Fix $E>0$.
If $h\in H^1(\T)$ and $g\in L^2(\T)$ satisfy
\begin{equation}\label{eq:Cauchy-test}
 \ip g{h_a}-\ip h{g_a}=0\qquad\hbox{for every }a,
\end{equation}
then they are the Cauchy data of a unique interior Helmholtz solution
$u\in H^{3/2}(\Omega)$. Also $\OO_E$ is injective and
\begin{equation}\label{eq:fixed}
 \dim\ker(V_E-I)\le\dim\ker(-\Delta_\Omega^N-E).
\end{equation}
If $\mu>0$ has Neumann multiplicity $m$ and $h_1,\ldots,h_m$ are the
boundary traces of an orthonormal basis $\phi_1,\ldots,\phi_m$ of that eigenspace, the Herglotz
normal traces at $\mu$ are dense in
\begin{equation}\label{eq:compatible}
 X_0=\{g\in X:\ip{h_j}g=0\ (1\le j\le m)\}.
\end{equation}
At a nonresonant energy they are dense in $X$.
\end{lemma}
\begin{proof}
The outgoing single-layer potential of density $g/|\gamma'|$ minus the
double-layer potential of density $h$ has zero far field by
\eqref{eq:Cauchy-test}, after conjugation and reversal of the plane-wave
direction. It vanishes in the connected exterior. The jump relations
give an interior $H^1$ solution with the prescribed data. Smooth-boundary
regularity upgrades it to $H^{3/2}$. Two such solutions have zero Cauchy
data and hence agree.

If $f=W_Ev$ and $f_0=0$, its first row gives
$0=f_0'=-(ik/\sqrt2)\gamma'f_1$, and successive rows give every $f_n=0$.
This proves injectivity of $\OO_E$. For $V_Ev=v$, pair $f$ with a driven
Herglotz vector and integrate. One gets $\ip{f_0}{g_a}=0$ for every $a$.
The negative of the double-layer potential of $f_0$ is zero outside
and has interior trace $f_0$ and homogeneous Neumann data. This
maps the fixed space injectively into the Neumann eigenspace and proves
\eqref{eq:fixed}.

For density in $X$, an annihilator is a trace $h\in X^*=H^{1/2}$.
The same double-layer argument identifies it with the trace of a Neumann
eigenfunction. Conversely each such trace annihilates every $g_a$ by
Green's identity. The Hahn--Banach theorem proves the exact closure
\eqref{eq:compatible}. The $h_j$ are independent by zero-Cauchy-data
uniqueness. This proof uses the dual topology $H^{-1/2}$, not merely an
$L^2$ approximation assertion.
\end{proof}

At a Neumann-nonresonant energy let $\NN(E):X\to X^*$ be the map
$g\mapsto u\circ\gamma$ for $(-\Delta-E)u=0$ and
$|\gamma'|\partial_\nu u\circ\gamma=g$. By \eqref{eq:fixed} the Cayley
transform
\[
 \KK_E=i(I+V_E)(I-V_E)^{-1},\qquad
 \Dom\KK_E=\ran(I-V_E)
\]
is densely defined and self-adjoint. Its resolvents are
$(\KK_E+i)^{-1}=(I-V_E)/(2i)$ and
$(\KK_E-i)^{-1}=(I-V_E)V_E^*/(2i)$, and are compact.
Set $\Aop_E=2\NN(E)+i(\TT_E-\TT_E^*)$ on $L^2$.
Endpoint algebra in \eqref{eq:endpoint} gives
\begin{equation}\label{eq:factor}
 \Aop_Eg_a=-\OO_E\KK_E\OO_E^*g_a,
 \qquad\{\OO_E^*g_a:a\in L^2\}=\Dom\KK_E.
\end{equation}
Indeed the endpoint identity first places $\OO_E^*g_a$ in the domain;
then $i(I-V_E)^{-1}V_E=\KK_E/2-iI/2$ and
\eqref{eq:OTadj} give the formula. Both assertions are on the stated domain of $\KK_E$; the first
is not an extension of an unbounded composition by density. An eigenvalue $e^{i\theta}$ of $U_E$,
$0<\theta<2\pi$, corresponds to $\cot(\theta/2)$ for $\KK_E$.

\section{Antiholomorphic transmutation}\label{sec:transmutation}

Let $\mathbb D=\{z\in\mathbb C:|z|<1\}$ and choose a Riemann
map $F:\mathbb D\to\Omega$ with $F(1)=0$. We use conformal
boundary coordinates $\gamma(\theta)=F(e^{i\theta})$. The map
$F$ is a smooth diffeomorphism on the closures
and $F'$ has no boundary zero \cite{Pommerenke}. Reparametrization does
not change $V_E$. Let $\mathscr A^s$ denote antiholomorphic functions on
$\Omega$ whose pulled-back traces belong to $H^s(\T)$; their traces are
exactly the Fourier space with modes $n\le0$.

For $h\in\mathscr A^s$, $s\ge1/2$, its antiholomorphic primitive $Jh$
is normalized by $(Jh)(0)=0$. It is single valued because $\Omega$ is
simply connected. The identity
$(Jh\circ\gamma)'=\overline{\gamma'}(h\circ\gamma)$ implies
\begin{equation}\label{eq:J-bound}
 \norm{Jh\circ\gamma}_{H^{s+1}}\le C_s\norm{h\circ\gamma}_{H^s}.
\end{equation}
Evaluation is used only on the primitive, which belongs to $H^{s+1}$.
Define
\begin{equation}\label{eq:Vekua}
 \VV_Eh(z)=\sum_{j=0}^\infty\frac{(-E/4)^j}{j!}z^jJ^jh(z),
 \qquad \JJ_Eh=|\gamma'|\partial_\nu\VV_Eh\circ\gamma.
\end{equation}
The harmonic input in $\JJ_E$ is identified with its boundary trace.
This is the antiholomorphic version of Vekua transmutation
\cite{Vekua,MHP}. The construction below uses single-valued primitives
and \eqref{eq:J-bound}; no star-shapedness assumption is needed.

\begin{lemma}\label{lem:transmutation}
The series \eqref{eq:Vekua} is entire in $E$ in the indicated Sobolev
operator norms and satisfies $(-\Delta-E)\VV_Eh=0$. The map $\VV_E$ is
injective. For $s\ge1/2$, locally uniformly in $E$,
\begin{equation}\label{eq:gains}
 (\VV_E-I):\mathscr A^s\longrightarrow H^{s+1}(\T),
 \qquad (\JJ_E-|D|):\mathscr A^s\longrightarrow H^s(\T)
\end{equation}
are bounded, with norms $O(E)$ at zero. Here the first map in
\eqref{eq:gains} means the boundary trace of $\VV_E-I$.
In the interior one also has
\begin{equation}\label{eq:interior-transmutation}
 \VV_E:\mathscr A^s\longrightarrow H^{s+1/2}(\Omega),\qquad
 \VV_E-I:\mathscr A^s\longrightarrow H^{s+3/2}(\Omega).
\end{equation}
The second map has norm $O(E)$ near zero. The conormal map $\JJ_E$
is bounded $\mathscr A^s\to H^{s-1}(\T)$, with its weak conormal
interpretation when $s=1/2$.
For $h\in\mathscr A^1$, if $h_0=h(0)$, then
\begin{equation}\label{eq:anti-endpoint}
 \OO_E^*\JJ_Eh=i(V_E^*-I)e_0h_0,
 \qquad (\VV_Eh)\circ\gamma=h_0\OO_Ee_0-i\TT_E\JJ_Eh.
\end{equation}
\end{lemma}
\begin{proof}
On the trace space $H^s$, the primitive is bounded and gains one
derivative by \eqref{eq:J-bound}. Iterating its $H^s$ bound, and
using the gain in the last application, bounds $J^j$ from $H^s$
to $H^{s+1}$ by $C_s^j$ for $j\ge1$. Multiplication by $\gamma^j$ has norm at
most $C(1+j)^rR^j$ for a fixed integer $r>s+1$. The factorial thus
proves convergence, the first gain in \eqref{eq:gains}, and parameter
continuity. Harmonic extension takes an antiholomorphic trace in
$H^r(\T)$ into $H^{r+1/2}(\Omega)$. Apply this to $h$ and to each
$J^jh$, $j\ge1$, using one primitive gain before summing. Multiplication
by $z^j$ on the fixed smooth domain has a bound of the form
$C(1+j)^r R_0^j$ at every fixed Sobolev order. The factorial gives
\eqref{eq:interior-transmutation}, including convergence in the interior.
Termwise differentiation gives
\[
 \partial\bar\partial(z^jJ^jh)=jz^{j-1}J^{j-1}h,
\]
which proves the equation. For $j\ge1$ the conormal trace of a summand is
\[
 -ij\gamma'\gamma^{j-1}(J^jh)\circ\gamma
 +i\overline{\gamma'}\gamma^j(J^{j-1}h)\circ\gamma.
\]
Its summation proves the conormal claims. Harmonic conormal data in these
coordinates are $|D|(h\circ\gamma)$.

For injectivity, termwise differentiation gives
$\partial^n\VV_Eh=(-E/4)^n\VV_EJ^nh$. The same bounds justify the double
series in the identity
\begin{equation}\label{eq:V-inverse}
 h=\sum_{n=0}^\infty\frac{(-z)^n}{n!}\partial^n(\VV_Eh).
\end{equation}
For each total degree $\ell$, the inner coefficient is
$\sum_{n=0}^\ell(-1)^n/(n!(\ell-n)!)$, zero unless $\ell=0$.
This proves \eqref{eq:V-inverse}, including convergence in boundary $H^s$.

Antiholomorphic polynomials are dense in $\mathscr A^1$ in its trace
norm. To see the stated strength of this approximation, first use radial
smoothing in conformal coordinates. For an input smooth on the closure,
approximate the derivative of its conjugate uniformly by polynomials
using Mergelyan's theorem \cite{Rudin}; integrate and fix the value at the
boundary origin. Paths in the closure can be chosen of uniformly bounded
length. This gives approximation of both boundary values and tangential
derivatives, hence in $H^1$. For $h=\bar z^m$, $m\ge0$, the series is
\[
 \sum_{j\ge0}\frac{(-E/4)^jm!}{j!(m+j)!}z^j\bar z^{m+j}.
\]
The Herglotz wave with density $e^{-im\phi}$ is
$(-ik/2)^m/m!$ times this series. Its initial vector is zero for $m>0$, and
$e_0/\sqrt2$ for $m=0$. The identities \eqref{eq:anti-endpoint} therefore
follow from \eqref{eq:trace-driven}--\eqref{eq:endpoint} for polynomials,
and then for $\mathscr A^1$ by the proved bounds. The endpoint value in these identities is only used for $H^1$
input traces.
\end{proof}

The reverse kernel inclusion uses the following consequence of
standard elliptic trace theory \cite{LM,McLean}. Normal derivatives are
weak conormal traces; this distinction matters at the Sobolev index
$3/2$.

\begin{lemma}\label{lem:generalized-traces}
Fix $E>0$ and a smooth domain. An $L^2$ distributional solution of
$(-\Delta-E)w=0$ has a generalized Dirichlet trace. If that trace belongs
to $H^1(\partial\Omega)$, then $w\in H^{3/2}(\Omega)$ and its weak
normal derivative belongs to $L^2(\partial\Omega)$. Such solutions satisfy
\begin{equation}\label{eq:Cauchy-estimate}
 \norm w_{H^{3/2}(\Omega)}\le C(\Omega,E)
 \bigl(\norm{\operatorname{Tr}w}_{H^1}
       +\norm{\partial_\nu^\gamma w}_{L^2}\bigr).
\end{equation}
Writing $h=\operatorname{Tr}w$ and $g=\partial_\nu^\gamma w$, the
traces of the $L^2$ Helmholtz solutions $\partial w$ and $\bar\partial w$
obey, in distributions on the boundary,
\begin{equation}\label{eq:generalized-identities}
 \begin{split}
 h'&=\gamma'\operatorname{Tr}(\partial w)
          +\overline{\gamma'}\operatorname{Tr}(\bar\partial w),\\
 g&=-i\gamma'\operatorname{Tr}(\partial w)
          +i\overline{\gamma'}\operatorname{Tr}(\bar\partial w).
 \end{split}
\end{equation}
\end{lemma}
\begin{proof}
The generalized Dirichlet trace of an $L^2$ solution is defined by
Green's identity in $H^{-1/2}(\partial\Omega)$; see \cite{LM,McLean}.
Let $w_b$ be the harmonic extension of its $H^1$ trace. Then
$w_b\in H^{3/2}$, whereas $w-w_b$ has zero generalized trace and
$\Delta(w-w_b)=-Ew\in L^2$. The regularity of the zero-Dirichlet
Poisson problem gives $w-w_b\in H^2\cap H_0^1$. Consequently
\[
 \norm w_{H^{3/2}}\le C(\Omega,E)
       (\norm{\operatorname{Tr}w}_{H^1}+\norm w_{L^2}),
\]
and the normal derivative is in $L^2$ by the harmonic
Dirichlet-to-Neumann map and the $H^2$ correction. The additional
$L^2$ term can be removed when both Cauchy norms are included:
otherwise a sequence normalized in $H^{3/2}$, compact in $L^2$,
would converge to a nonzero solution with zero Cauchy data.
Uniqueness for those data is a contradiction. This proves
\eqref{eq:Cauchy-estimate}, also at eigenvalues.
The tangential and normal identities \eqref{eq:generalized-identities}
follow from Green's identity, first for smooth functions and then for
these generalized traces. In particular
$\operatorname{Tr}(\partial w)=(h'+ig)/(2\gamma')$.
\end{proof}

We next identify the kernel of the observation adjoint.
Lemma~\ref{lem:transmutation} gives one inclusion; the reverse
inclusion requires control of all holomorphic derivatives.

\begin{lemma}\label{lem:exact-kernel}
For every $E>0$, including a Neumann resonance,
\begin{equation}\label{eq:exact-kernel}
 \ker\OO_E^*=
 \{\JJ_Eh:h\in\mathscr A^1,\ h(0)=0\}\quad\hbox{in }L^2(\T).
\end{equation}
\end{lemma}
\begin{proof}
One inclusion follows from \eqref{eq:anti-endpoint}. For the converse
let $g\in L^2$ with $\OO_E^*g=0$ and solve the driven equation
\eqref{eq:driven} with $y(0)=0$. The variation formula gives
$y(L)=0$, so $y\in H^1(\T;\HH)$. Set $f_0=\sqrt2 y_0$ and $f_n=y_n$
for $n\ge1$. Pairing with any driven Herglotz vector gives
\[
 \frac d{d\theta}\ip y{y_a}=\frac i2
       (\overline g h_a-\overline{f_0}g_a).
\]
Both vectors are periodic. Lemma~\ref{lem:Cauchy} therefore gives a
Helmholtz solution $u$ with trace $f_0$ and conormal trace $g$.

We reconstruct the derivatives without presupposing their regularity.
Put $u_0=u$. The first driven row gives
$f_0'+ig=-ik\gamma'f_1$. By \eqref{eq:generalized-identities},
$u_1=(2i/k)\partial u$ has generalized trace $f_1\in H^1$;
Lemma~\ref{lem:generalized-traces} therefore makes $u_1$ an
$H^{3/2}$ solution. Inductively, suppose
$u_n=(2i/k)^n\partial^n u\in H^{3/2}$ with trace $f_n$.
For $n\ge1$ the equation gives
$\bar\partial u_n=-ik u_{n-1}/2$. The tangential identity and the
$n$th driven row yield
\[
 \gamma'\operatorname{Tr}(\partial u_n)
   =f_n'+\frac{ik}{2}\overline{\gamma'}f_{n-1}
   =-\frac{ik}{2}\gamma'f_{n+1}.
\]
Thus $u_{n+1}$ has generalized trace $f_{n+1}\in H^1$ and is again
in $H^{3/2}$. The normal identity gives, for $n\ge1$,
\[
 \operatorname{Tr}u_n=f_n,\qquad
 \partial_\nu^\gamma u_n=\frac k2
        (\overline{\gamma'}f_{n-1}-\gamma'f_{n+1}).
\]
The Cauchy estimate is for the same equation at every $n$, so its
constant does not depend on $n$. Summing its square and using
$y\in H^1(\T;\HH)$ gives
\begin{equation}\label{eq:derivative-bound}
 \sum_{n\ge0}\norm{u_n}_{H^{3/2}(\Omega)}^2
 \le C(\Omega,E)\bigl(\norm y_{H^1(\T;\HH)}^2+\norm g_{L^2}^2\bigr)<\infty.
\end{equation}
In particular the norms of $\partial^n u$ grow at most geometrically.

Since multiplication by $z^\ell$ on $H^{3/2}(\Omega)$ is bounded
by a polynomial in $\ell$ times $R^\ell$, for a fixed $R$, the
factorials below dominate all multiplier and derivative bounds.
For $j\ge0$ define
\[
 H_j(z)=\sum_{\ell\ge0}\frac{(-z)^\ell}{\ell!}\partial^{j+\ell}u(z).
\]
The series converges in $H^{3/2}(\Omega)$ and in boundary $H^1$ by
\eqref{eq:derivative-bound}. Differentiation and cancellation show
$\partial H_j=0$ and, for $j\ge1$,
$\bar\partial H_j=-(E/4)H_{j-1}$. All $f_n(0)$ vanish, so
$H_j(0)=0$. Consequently $H_j=(-E/4)^jJ^jH_0$ and
$H_0\in\mathscr A^1$ has value zero at the origin. The absolutely
convergent binomial double series yields
\[
 u=\sum_{j\ge0}\frac{z^j}{j!}H_j=\VV_EH_0.
\]
Taking conormal traces proves the reverse inclusion in
\eqref{eq:exact-kernel}.
\end{proof}

\section{The periodic boundary form}\label{sec:periodic}

The boundary origin may be chosen separately on a neighborhood of a
fixed positive energy. This changes $V_E$ by unitary conjugation and
hence does not change its phases or determinant.

\begin{lemma}\label{lem:cut}
At every $E_0>0$ there is a choice of boundary origin and a neighborhood
of $E_0$ on which
$c_E=(V_E^*-I)e_0\ne0$. Put $s_E=(I+V_E^*)e_0$, $d_E=is_E$ and
\begin{equation}\label{eq:Bcut}
 B_E=\frac{d_E\ip{c_E}{\,\cdot\,}+c_E\ip{d_E}{\,\cdot\,}}{\norm{c_E}^2}
       -\frac{\ip{c_E}{d_E}}{\norm{c_E}^4}
                   c_E\ip{c_E}{\,\cdot\,},
 \qquad \Pi_E^c=I-\frac{c_E\ip{c_E}{\,\cdot\,}}{\norm{c_E}^2}.
\end{equation}
Here $\Pi_E^c$ is the orthogonal projection onto $c_E^\perp$.
The operator $B_E$ is self-adjoint of rank at most two,
$B_Ec_E=is_E$, and both families are norm-continuous. At nonresonant energies define
\begin{equation}\label{eq:periodic-operators}
 \PP_E=\Aop_E+\OO_E B_E\OO_E^*,\qquad
 T_E^\sharp=\Pi_E^c\OO_E^*.
\end{equation}
The first operator extends to a bounded Hermitian map $X\to X^*$,
in the sense that
$\ip f{\PP_Eg}=\overline{\ip g{\PP_Ef}}$ for $f,g\in X$.
The second extends compactly from $X$ to $\HH$. Moreover
\begin{equation}\label{eq:periodic-null}
 \PP_E\JJ_E=0,\qquad T_E^\sharp\JJ_E=0
       \quad\hbox{on }\mathscr A^{1/2},
\end{equation}
and, on $L^2$ at every energy in the neighborhood,
\begin{equation}\label{eq:projected-kernel}
 \ker T_E^\sharp=\ran(\JJ_E:\mathscr A^1\to L^2).
\end{equation}
The map $\JJ_E:\mathscr A^{1/2}\to X$ is injective there, including at a
resonant energy.
\end{lemma}
\begin{proof}
If every rebasing at $\gamma(t)$ fixed $e_0$, the original $V_{E_0}$
would fix $W_{E_0}(t)^*e_0$ for every $t$. These vectors span a dense
space by injectivity of $\OO_{E_0}$. Thus $V_{E_0}=I$, contradicting
\eqref{eq:fixed}. Choose a rebasing for which $c_{E_0}\ne0$ and use
continuity. Since
$\ip{c_E}{d_E}=-2\operatorname{Im}\ip{e_0}{V_Ee_0}\in\mathbb R$,
\eqref{eq:Bcut} is self-adjoint and maps $c_E$ to $is_E$.

For $v\in\HH$ the endpoint jump of $\OO_Ev$ is $\ip{c_E}v$.
For $g\in L^2$, the jump of $\Aop_Eg$ is $i\ip{s_E}{\OO_E^*g}$.
The added operator in \eqref{eq:periodic-operators} cancels this jump
because $B_Ec_E=is_E$ and the inner product is conjugate-linear in its
first variable. The bounded Sobolev extension for $\PP_E$ is proved
quantitatively in Lemma~\ref{lem:principal} below.
Also $\OO_E\Pi_E^c:\HH\to H^1(\T)$ is Hilbert--Schmidt: it is a
bounded smooth interval observation with matching endpoint values, and
its first derivative has a rank-one row of uniformly bounded norm.
By duality and the compact embedding $H^{-1/2}\hookrightarrow H^{-1}$,
$T_E^\sharp:X\to\HH$ is compact and norm-continuous.

Let $p_E=\VV_E1=J_0(k|z|)$ and $q_E=\JJ_E1$, where $J_0$ denotes
the Bessel function of the first kind of order zero. Formula
\eqref{eq:anti-endpoint} and endpoint algebra give
\begin{equation}\label{eq:radial}
 \OO_E^*q_E=ic_E,\qquad \Aop_Eq_E=\OO_Es_E.
\end{equation}
The second equality is for nonresonant energies. Thus $\PP_Eq_E=0$.
For an antiholomorphic $h$ with value zero at the origin,
\eqref{eq:anti-endpoint} gives $\OO_E^*\JJ_Eh=0$ and
$\Aop_E\JJ_Eh=0$. One can obtain the latter either from polynomial
Herglotz approximation and \eqref{eq:factor}, or directly from
\eqref{eq:anti-endpoint} and $\TT_E+\TT_E^*=\OO_E\OO_E^*$.
This proves \eqref{eq:periodic-null} first on $\mathscr A^1$, then on
$\mathscr A^{1/2}$ by density and the proved Sobolev bounds.

If $T_E^\sharp g=0$ for $g\in L^2$, its observation is a multiple of
$c_E$. Subtract a suitable multiple of $q_E$, whose observation is
$ic_E$, and apply Lemma~\ref{lem:exact-kernel}. This proves
\eqref{eq:projected-kernel}, also at a resonance.

Finally suppose $\JJ_Eh=0$ for $h\in\mathscr A^{1/2}$.
By \eqref{eq:interior-transmutation}, $u=\VV_Eh$ belongs to $H^1$.
It is a weak Neumann eigenfunction, or zero at nonresonance, and is smooth up to
the boundary. Since $(\VV_E-I)h$ has $H^{3/2}$ trace, $h$ belongs to
$\mathscr A^{3/2}$, so its point value is legitimate. The first identity
in \eqref{eq:anti-endpoint} gives $ic_Eh(0)=0$, and hence $h(0)=0$.
The second gives $\operatorname{Tr}u=0$. Uniqueness for zero Cauchy data
makes $u=0$, and injectivity of $\VV_E$ makes $h=0$.
\end{proof}

We give the regularity details of the periodic correction. They will
also permit a bootstrap from the dual trace space back to $L^2$.

\begin{lemma}\label{lem:principal}
At a nonresonant energy the normalized self-adjoint operator
\begin{equation}\label{eq:normalized}
 \widehat P_E=\Lambda^{1/2}\PP_E\Lambda^{1/2}
       =4\Pi_++C_E
\end{equation}
is bounded on $L^2$, where $C_E$ is compact. Locally its compact part
is norm-continuous. For any fixed $0<\delta<1/2$, the compact remainder
$C_E$, not $\widehat P_E$, is bounded $H^t\to H^{t+\delta}$ for
$0\le t\le1/2$, locally uniformly.
If $E_0=\mu$ is a Neumann eigenvalue, then
\begin{equation}\label{eq:pole}
 \widehat P_E=-\frac2{E-\mu}\sum_{j=1}^m
    (\Lambda^{1/2}h_j)\ip{\Lambda^{1/2}h_j}{\,\cdot\,}
       +\widehat P_{\mathrm{reg}}(E),
 \quad \widehat P_{\mathrm{reg}}(E)=4\Pi_++C_{\mathrm{reg}}(E),
\end{equation}
where $C_{\mathrm{reg}}(E)$ has the same continuity and regularizing
properties across $\mu$.
Moreover
\begin{equation}\label{eq:Jhat}
 \widehat J_E:=\Lambda^{-1/2}\JJ_E\Lambda^{-1/2}
 :\mathcal H_{\le0}\longrightarrow L^2
\end{equation}
is the inclusion plus a compact operator gaining one derivative.
\end{lemma}
\begin{proof}
On mean-zero conormal data the harmonic Neumann map, up to its additive
constant, is $|D|^{-1}$. Choose the harmonic solution $u_0$ with zero
volume mean. Then
\begin{equation}\label{eq:harmonic-comparison}
 u_E=u_0+E(A_N-E)^{-1}u_0,
\end{equation}
where $A_N$ is the Neumann Laplacian. The correction maps
$H^{-1/2}$ boundary data into $H^3(\Omega)$, and hence its trace into
$H^{5/2}$. Constants contribute smooth finite-rank operators. Resolvent
expansion at $\mu$, or the same formula using a fixed nonresonant
reference energy for arbitrary data, gives
\begin{equation}\label{eq:N-pole}
 \NN(E)=-\frac1{E-\mu}\sum_{j=1}^m h_j\ip{h_j}{\,\cdot\,}
                       +\NN_{\mathrm{reg}}(E).
\end{equation}
For a direct normalization check, let $P_{-1}g$ solve the Neumann problem
at energy $-1$. Then
$u_E=P_{-1}g+(E+1)(A_N-E)^{-1}P_{-1}g$ and
$\ip{\phi_j}{P_{-1}g}=\ip{h_j}g/(\mu+1)$.
This proves \eqref{eq:N-pole}. Its regular remainder has the harmonic
principal part just described and the additional smoothness used here.

The periodic integral kernel of $2D^{-1}$, defined to vanish on constants,
is $i\operatorname{sgn}(\theta-s)-i(\theta-s)/\pi$ on the cut square $[0,L]^2$.
Subtracting it from $i(\TT_E-\TT_E^*)+\OO_EB_E\OO_E^*$ leaves the kernel
\begin{align}\label{eq:regular-kernel}
 r_E(\theta,s)={}&i\operatorname{sgn}(\theta-s)(w_E(\theta,s)-1)
                  +i(\theta-s)/\pi\notag\\
 &+\ip{e_0}{W_E(\theta)B_EW_E(s)^*e_0},
 \qquad w_E(\theta,s)=\ip{e_0}{W_E(\theta)W_E(s)^*e_0}.
\end{align}
Here $w_E(s,s)=1$ and $\partial_\theta w_E(s,s)=0$. Thus the first term
is $C^1$ across the diagonal. The endpoint cancellation in
Lemma~\ref{lem:cut} shows that opposite-edge values of \eqref{eq:regular-kernel}
match. The kernel is smooth on the pieces of the cut square and has
matching values, so it is in $H^{1+\delta}(\T^2)$ for every $\delta<1/2$.
One way to check this last assertion is to differentiate once: the
resulting bounded piecewise smooth functions, with jumps only across
finitely many smooth lines, belong to $H^\delta$ for $\delta<1/2$ by
the fractional difference-quotient norm. First derivatives need not
match at the cut. All these estimates are locally uniform and continuous.

If $r_{mn}$ are its Fourier matrix entries, then for $0\le t\le1/2$
\[
 \sum_{m,n}(1+|m|)^{1+2t+2\delta}(1+|n|)^{1-2t}|r_{mn}|^2
       \le C_\delta\norm{r_E}_{H^{1+\delta}(\T^2)}^2.
\]
Thus sandwiching the remainder by $\Lambda^{1/2}$ gives a compact
operator $H^t\to H^{t+\delta}$. The combined principal part is
$2|D|^{-1}+2D^{-1}=4D_+^{-1}\Pi_+$; its normalization is $4\Pi_+$ plus
an operator gaining one derivative. The Neumann comparison above proves
\eqref{eq:normalized} and \eqref{eq:pole}. At higher Sobolev orders
in the stated range, the same comparison uses the corresponding
smooth-boundary elliptic estimates. The pole-subtracted resolvent
and the kernels depend continuously on $E$ in these norms. This also proves the
$X\to X^*$ extension asserted in Lemma~\ref{lem:cut}.

Finally the principal part of \eqref{eq:Jhat} is
$\Lambda^{-1/2}|D|\Lambda^{-1/2}$ on nonpositive modes. It differs from
the inclusion by an operator gaining one derivative, including a
one-dimensional discrepancy on constants. Equation~\eqref{eq:gains}
gives the same gain for its remainder.
\end{proof}

\section{Uniform bounds for the Cayley transform}\label{sec:uniform}

We now fix any positive $E_0$, choose the local cut from
Lemma~\ref{lem:cut}, and work on its energy neighborhood. All neighborhoods
below may be decreased. They contain no other Neumann eigenvalue.
The symbol $m$ is zero at a nonresonant $E_0$, and otherwise its Neumann
multiplicity. When $m=0$, regular parts mean the original operators,
$X_0=X$, and $Z=\mathcal H_+$. Empty sums and compatibility
conditions are then omitted.

\begin{lemma}\label{lem:chart}
There is a fixed map $R:\mathcal H_+\to L^2$, equal to the inclusion plus
a finite-rank map with finite Fourier input and output vectors, such that
\begin{equation}\label{eq:chart}
 \mathfrak B_E(a,z)=\widehat J_Ea+Rz:
 \mathcal H_{\le0}\oplus\mathcal H_+\longrightarrow L^2
\end{equation}
is invertible locally in $E$, with norm-continuous inverse. Put
\[
 \mathscr R_E=R^*\widehat P_ER,
 \qquad \mathscr S_E=T_E^\sharp\Lambda^{1/2}R.
\]
After its pole is subtracted, the first operator is $4I$ plus a
compact remainder having the regularizing properties of
Lemma~\ref{lem:principal}. The second operator is compact and
norm-continuous. At a resonance put
$\eta_j=R^*\Lambda^{1/2}h_j$ and
\[
 Z=\{z\in\mathcal H_+:\ip{\eta_j}z=0\ (1\le j\le m)\}.
\]
The space $Z$ has codimension $m$, its orthogonal complement consists
of smooth functions, and the pole of $\mathscr R_E$ vanishes on $Z$.
We write $\Pi_Z$ for the orthogonal projection onto $Z$.
\end{lemma}
\begin{proof}
The map $\widehat J_{E_0}\Pi_{\le0}+\Pi_+$ is $I+$ compact, hence Fredholm
of index zero. Its restriction to $\mathcal H_{\le0}$ is injective by
Lemma~\ref{lem:cut}. Consequently projection of its kernel onto the
positive input space is injective. A finite-rank map acting only on
positive inputs can therefore pair this kernel isomorphically with the
cokernel and make the whole operator invertible. Approximate its finitely
many input and output vectors by finite Fourier sums; invertibility is
open, so the correction can have the stated smooth vectors. Fix it as
$R-\Pi_+$. Continuity proves \eqref{eq:chart} locally.
The reduced operator assertions follow from Lemma~\ref{lem:principal}.

Every $\JJ_{E_0}h$ is the normal trace of a Helmholtz solution at $E_0$,
and hence is orthogonal to each $h_j$. Thus
$\Lambda^{1/2}h_j$ annihilates $\ran\widehat J_{E_0}$. If a combination
of the $\eta_j$ is zero, the corresponding combination of
$\Lambda^{1/2}h_j$ also annihilates $\ran R$, and hence all of $L^2$ by
\eqref{eq:chart}. The independence in Lemma~\ref{lem:Cauchy} gives the
claim. Formula~\eqref{eq:pole} now gives
\begin{equation}\label{eq:reduced-pole}
 \mathscr R_E=-\frac2{E-E_0}\sum_{j=1}^m
                      \eta_j\ip{\eta_j}{\,\cdot\,}
                         +\mathscr R_{\mathrm{reg}}(E),
 \quad \mathscr R_{\mathrm{reg}}(E)=4I+\text{compact}.
\end{equation}
The fact that $R$ is fixed, rather than energy dependent, makes $Z$ fixed
and cancels the pole there exactly.
\end{proof}

\begin{lemma}\label{lem:injective-reduced}
The restriction $\mathscr S_{E_0}|_Z$ is injective. At a nonresonant energy
this means injectivity on all of $\mathcal H_+$.
\end{lemma}
\begin{proof}
Let $z\in Z$ satisfy $\mathscr S_{E_0}z=0$ and set
$g=\Lambda^{1/2}Rz\in X$. Initially this is only a dual trace, not an
$L^2$ function. It belongs to $X_0$ from \eqref{eq:compatible}.

For Herglotz data at $E_0$, endpoint identities without an inverse give
\begin{equation}\label{eq:physical-regular}
 \PP_{\mathrm{reg}}(E_0)g_a
   =\OO_{E_0}\xi_a-2\sum_{j=1}^m h_j\ip{\phi_j}{u_a},
 \quad
 \xi_a=\sqrt2\bigl[(I+V^*)+iB(V^*-I)\bigr]v_a\in c^\perp.
\end{equation}
Here $\PP_{\mathrm{reg}}$ has the Neumann resolvent pole removed, and all unmarked
operators on the right are at $E_0$. At nonresonance it is simply $\PP$.
To verify this, use
$2h_a+i(\TT-\TT^*)g_a=\sqrt2\OO(I+V^*)v_a$ and
$\OO^*g_a=\sqrt2 i(V^*-I)v_a$. The identity $Bc=is$ shows directly that
$\ip c{\xi_a}=0$. At resonance $\NN_{\mathrm{reg}}(E_0)g_a$ is the trace
of the solution orthogonal to the eigenspace, namely
$h_a-\sum h_j\ip{\phi_j}{u_a}$, which proves \eqref{eq:physical-regular}.

Pair \eqref{eq:physical-regular} with $g$. The finite sum vanishes by
compatibility. The observation term vanishes because
$T_{E_0}^\sharp g=0$ and $\xi_a\in c^\perp$.
These pairings are legitimate on $X$: $\OO\xi_a$ is periodic $H^1$,
whereas the raw observation need not be periodic. Density in
Lemma~\ref{lem:Cauchy} implies
\[
 \ip g{\PP_{\mathrm{reg}}(E_0)\widetilde g}=0
                       \qquad(\widetilde g\in X_0).
\]
In particular, with $\widetilde g=\Lambda^{1/2}Rw$, $w\in Z$, this says
\begin{equation}\label{eq:bootstrap}
 \Pi_Z\mathscr R_{\mathrm{reg}}(E_0)z=0,
 \qquad 4z=-\Pi_Z C z,
 \quad C=\mathscr R_{\mathrm{reg}}(E_0)-4I.
\end{equation}
The projection $\Pi_Z$ differs from the identity by a smooth finite-rank
map. Take $\delta=1/4$ in Lemma~\ref{lem:principal}; twice applying
\eqref{eq:bootstrap} gives $z\in H^{1/2}$, so $g\in L^2$.

We may therefore apply \eqref{eq:projected-kernel}. It gives
$g=\JJ_{E_0}h$ for $h\in\mathscr A^1$. After normalization this lies in
$\ran\widehat J_{E_0}$. But it is also $Rz$. Uniqueness in the decomposition
\eqref{eq:chart} forces $z=0$. This proves injectivity using the kernel identity only at its proved $L^2$ regularity.
\end{proof}

\begin{lemma}\label{lem:coercive}
Let $\mathcal Z$ and $\mathcal Y$ be Hilbert spaces. Let
$H=\alpha I+C$ be self-adjoint on $\mathcal Z$, with $\alpha>0$
and $C$ compact, and let $S:\mathcal Z\to\mathcal Y$ be bounded
and injective. There are constants $A>0$ and $\varepsilon>0$ such that
\[
 H+A S^*S\ge\varepsilon I.
\]
\end{lemma}
\begin{proof}
If no integer $n$ gives a strictly positive lower bound for
$H+nS^*S$, choose unit vectors $z_n$ with
$\ip{z_n}{Hz_n}+n\norm{Sz_n}^2\le1/n$.
The first term is bounded, so $Sz_n\to0$. Pass to a weakly convergent
subsequence. Its limit lies in $\ker S$ and is zero. Compactness gives
$\ip{z_n}{Cz_n}\to0$, so $\ip{z_n}{Hz_n}\to\alpha$, a contradiction.
The resulting inequality remains valid with $\varepsilon/2$ in
place of $\varepsilon$ under sufficiently small operator-norm
perturbations of $H$ and $S$, with the same $A$.
\end{proof}

\begin{proposition}\label{prop:rank-bound}
If $E_0>0$ is nonresonant, $\KK_E\le A I$ uniformly on a neighborhood of
$E_0$. If $E_0=\mu>0$ has Neumann multiplicity $m$, there are $A>0$ and
$a>0$ such that
\begin{equation}\label{eq:rank-bound}
 \operatorname{rank}\mathbf1_{(A,\infty)}(\KK_E)\le m
 \qquad(0<|E-\mu|<a).
\end{equation}
Only nonresonant energies are used in this formula.
\end{proposition}
\begin{proof}
Apply Lemma~\ref{lem:coercive} on the fixed space $Z$ to
$\Pi_Z\mathscr R_{\mathrm{reg}}(E_0)|_Z$ and
$\mathscr S_{E_0}|_Z$. The hypotheses follow from
Lemmas~\ref{lem:chart} and \ref{lem:injective-reduced}. By continuity and
\eqref{eq:reduced-pole}, locally and off resonance,
\begin{equation}\label{eq:reduced-bound}
 \ip z{\mathscr R_Ez}\ge -A_0\norm{\mathscr S_Ez}^2\qquad(z\in Z).
\end{equation}
For any $g\in X$ write
$\Lambda^{-1/2}g=\widehat J_Ea+Rz$ using \eqref{eq:chart}.
Equations \eqref{eq:periodic-null} imply
\[
 \ip g{\PP_Eg}=\ip z{\mathscr R_Ez},\qquad
 T_E^\sharp g=\mathscr S_Ez.
\]
Impose the at most $m$ linear conditions $z\in Z$ on $g$.
For $g\in L^2$ obeying them, \eqref{eq:reduced-bound} and boundedness of
$B_E$ give
\begin{equation}\label{eq:boundary-ineq}
 \ip g{\Aop_Eg}\ge
 -(A_0+\norm{B_E})\norm{\OO_E^*g}^2.
\end{equation}
The norm of $B_E$ is bounded on the chosen neighborhood.

Suppose $\KK_E$ had an $(m+1)$-dimensional spectral subspace strictly
above a larger uniform constant $A$. Lift a basis to Herglotz normal
traces by the range identity in \eqref{eq:factor}. Their span
maps isomorphically to that subspace. At least one nonzero linear
combination obeys the $m$ linear conditions, and \eqref{eq:factor} makes
its quadratic form strictly less than $-A\norm{\OO_E^*g}^2$, contradicting
\eqref{eq:boundary-ineq}. This proves \eqref{eq:rank-bound}.
For $m=0$ the argument gives the full upper bound. In particular it gives
a uniform spectral gap immediately above $1$ for $U_E$ off the
Neumann spectrum. It does not assume that every phase arriving at $1$
has a nonzero limiting eigenvector.
\end{proof}

\section{Small-energy normalization}\label{sec:small}

The counting argument requires a normalization at zero. The following
rank-one decomposition shows that the zero Neumann mode contributes
exactly one. For a unitary sufficiently close to the identity, a
principal eigenphase means an argument in $(-\pi,\pi)$.

\begin{lemma}\label{lem:initial}
For every smooth Jordan domain and all sufficiently small $E>0$,
\begin{equation}\label{eq:initial}
 \frac{E|\Omega|/2-\Phi_-(E)}{2\pi}=1.
\end{equation}
\end{lemma}
\begin{proof}
Restrict the energy interval to $(0,\mu_1(\Omega))$ and keep a fixed
conformal parametrization and base point. With $p_E=\VV_E1$, $q_E=\JJ_E1$, set
$b_E=\Aop_Eq_E$ and $a_E=\ip{q_E}{b_E}$. Equations
\eqref{eq:radial}, \eqref{eq:curvature} and \eqref{eq:det} give
\[
 b_E=\OO_E(I+V_E^*)e_0=2+\beta_E,\quad
 \norm{\beta_E}_{C^1[0,L]}=O(E),\quad
 a_E=2\operatorname{Im}\ip{e_0}{V_Ee_0}
           =-2E|\Omega|+o(E)<0.
\]
The zeroth-row matrix elements of the transport have no linear term in
$k$, which justifies the $C^1$ estimate. The jump of $\Aop_Eg$ is
$i\ip{b_E}g$, so the self-adjoint Schur complement
\begin{equation}\label{eq:small-Schur}
 C_E^{\rm sc}=\Aop_E-a_E^{-1}b_E\ip{b_E}{\,\cdot\,}
\end{equation}
has matching endpoint values and annihilates $q_E$ and every
$\JJ_Eh$ with $h(0)=0$.

On $\mathcal H_+$ put $S_+=D_+^{-1/2}$. Its positive Fourier block is
\begin{equation}\label{eq:small-positive}
 \ip g{C_E^{\rm sc}g}
      =\ip{S_+g}{(4I+K_E^{\rm sc})S_+g},\qquad
 \norm{K_E^{\rm sc}}=O(E).
\end{equation}
Here is an explicit check of the cut term in this assertion. The
remainder relative to $i(\TT_0-\TT_0^*)$, where
$\TT_0g(\theta)=\int_0^\theta g$, is
$r_E^0=i\operatorname{sgn}(\theta-s)(w_E-1)$, a $C^1$ kernel of size
$O(E)$ on the cut square. Direct endpoint evaluation gives
\[
 r_E^0(L,s)-r_E^0(0,s)=i\overline{\beta_E(s)},\qquad
 \beta_E(L)-\beta_E(0)=ia_E.
\]
Thus $r_E^0-\beta_E(\theta)\overline{\beta_E(s)}/a_E$ is a periodic
Hermitian $H^1(\T^2)$ kernel with norm $O(E)$. Its normalization by
$D_+^{1/2}$ on both sides is Hilbert--Schmidt, because
$mn\le(m^2+n^2)/2$. On mean-zero data the compression of the rank-one
term using $b_E$ equals the one using $\beta_E$. Formula
\eqref{eq:harmonic-comparison}, with zero volume mean, has norm $O(E)$
at zero after its smoothing gain, and direct Fourier integration gives
$\ip g{i(\TT_0-\TT_0^*)g}=2\sum_{n>0}|g_n|^2/n$.
The harmonic Neumann term contributes the other $2\sum|g_n|^2/n$.
These observations prove \eqref{eq:small-positive}.

For $g\in\mathcal H_-$ define $h_g\in\mathscr A^1$ by
\[
 h_g\circ\gamma=\sum_{n<0}\frac{g_n}{|n|}
                         (\ph_n-\ph_n(0)).
\]
It has value zero at the origin and harmonic conormal density $g$.
Equation~\eqref{eq:gains} gives $G_Eg:=\JJ_Eh_g=g+O(E)g$ in $L^2$.
Also
\[
 \int_0^L q_E\,d\theta=-E\int_\Omega p_E
                         =-E|\Omega|+O(E^2),\qquad
 q_E/E\longrightarrow-\tfrac12\operatorname{Im}(\overline\gamma\gamma')
                         \quad\hbox{in }L^2.
\]
The vector $\zeta_E=q_E/\ip{\ph_0}{q_E}$ therefore has constant Fourier
coefficient one and converges in $L^2$. The map
\[
 \widetilde B_E=\Pi_++G_E\Pi_-+\zeta_E\ip{\ph_0}{\,\cdot\,}
\]
is invertible for small $E$: its limit is $I$ plus the square-zero
rank-one map $(\zeta_0-\ph_0)\ip{\ph_0}{\,\cdot\,}$.
The operator $C_E^{\rm sc}$ annihilates $G_E\mathcal H_-$ and
$\zeta_E$. These inclusions and the invertibility of $\widetilde B_E$,
together with \eqref{eq:small-positive}, give $C_E^{\rm sc}\ge0$ on
all of $L^2$. No converse nullspace assertion is needed here.
Consequently $\indm\Aop_E=1$: it is at most one by \eqref{eq:small-Schur},
and at least one since $a_E<0$.

By \eqref{eq:factor}, every positive Cayley spectral subspace lifts to
a negative subspace for $\Aop_E$, so there is at most one positive
principal phase when $\norm{U_E-I}<1$. The principal logarithm has trace
$iE|\Omega|/2$ by \eqref{eq:det} and continuity down to zero. This positive
sum shows that at least one positive principal phase exists. Passing to
the negative argument branch subtracts $2\pi$ exactly once. This proves
\eqref{eq:initial}.
\end{proof}

\section{The counting inequality}\label{sec:counting}

Let $\Argm z\in(-2\pi,0)$ denote the negative argument of a
point $z$ on the unit circle other than $1$.
At any positive nonresonant $E$, Proposition~\ref{prop:rank-bound} with
$m=0$ shows that $\KK_E$ is bounded above. Its compact resolvent implies
that it has only finitely many positive eigenvalues. Hence $U_E$ has
only finitely many upper-semicircle eigenvalues, while
$|\Argm z|\le C|z-1|$ on the lower semicircle. The series
\[
 \Phi_-(E)=\sum_{z\in\sigma_p(U_E)\setminus\{1\}}\Argm z
\]
is therefore absolutely convergent, with multiplicities. The determinant
formula makes \eqref{eq:integer-intro} an integer. For $|z|=1$, $z\ne1$,
\[
 -\Argm z\ge |z-1|.
\]
The compact normal operator $U_E-I$ has its norm attained at a
nonzero spectral value whenever it is nonzero. Therefore
\begin{equation}\label{eq:phase-remainder}
 -\Phi_-(E)\ge\norm{U_E-I},\qquad
 M(E)\ge\frac{E|\Omega|}{4\pi}+\frac{\norm{U_E-I}}{2\pi}.
\end{equation}

On a compact nonresonant interval, the uniform upper Cayley bound gives
a uniform empty arc immediately above $1$. Put a logarithm cut in that
arc and normalize the logarithm to be zero at $1$. For this logarithm, $\log z=(z-1)f(z)$ with $f$ holomorphic near
the spectrum. Functional calculus and trace-norm continuity of
$U_E-I$ therefore make $\log U_E$ trace-norm continuous. Its trace
is $i\Phi_-(E)$. Thus $M$ is constant on each nonresonant
interval.

At a positive Neumann eigenvalue $\mu$ of multiplicity $m$, take $A,a$
from \eqref{eq:rank-bound}. Choose a sufficiently small $\varepsilon>0$
with $\cot(\varepsilon/2)>A$ and $e^{i\varepsilon}\notin\sigma(U_\mu)$.
This is possible since all spectral points other than $1$ are isolated
finite-multiplicity eigenvalues. By operator-norm continuity, the cut
through $e^{i\varepsilon}$ stays off the spectrum on a smaller neighborhood.
Let $\log_\varepsilon$ be the holomorphic branch on the plane cut
through $e^{i\varepsilon}$ with $\log_\varepsilon(1)=0$, and put
$\arg_\varepsilon U_E=-i\log_\varepsilon U_E$. Its values on the unit
circle have arguments in $(\varepsilon-2\pi,\varepsilon)$. Its trace is
continuous through $\mu$, even on actual fixed vectors. Therefore
\[
 J_\varepsilon(E)=\frac{E|\Omega|/2-
                 \tr\arg_\varepsilon U_E}{2\pi}
\]
is an integer constant through $\mu$. Write
$n_{(0,\varepsilon)}(U_E)$ for the number of eigenvalues of $U_E$
in the arc $\{e^{it}:0<t<\varepsilon\}$, counted with multiplicity.
At nonresonant nearby energies,
\begin{equation}\label{eq:jump-count}
 M(E)=J_\varepsilon(E)+n_{(0,\varepsilon)}(U_E),\qquad
 0\le n_{(0,\varepsilon)}(U_E)\le m.
\end{equation}
The second inequality is exactly \eqref{eq:rank-bound}. Thus the net
increase of $M$ across $\mu$ is at most $m$. Notice that no eigenvector
continuation at $1$ is used in this argument.

Lemma~\ref{lem:initial} gives $M=1$ below the first positive Neumann
eigenvalue. There are finitely many resonances below any fixed energy.
Summing \eqref{eq:jump-count} gives, for every nonresonant $E>0$,
\[
 M(E)\le 1+\sum_{0<\mu<E}\dim\ker(A_N-\mu)=N_N(E).
\]
Together with \eqref{eq:phase-remainder}, this yields
\begin{equation}\label{eq:quantitative-count}
 N_N(E)\ge\frac{E|\Omega|}{4\pi}+\frac{\norm{U_E-I}}{2\pi}.
\end{equation}
For a resonant energy $\mu>0$, approach from below through nonresonant
energies. The strict counting function is constant just below $\mu$
and equals $N_N(\mu)$ there. Norm continuity of $U_E$ therefore proves
\eqref{eq:quantitative-count} at $\mu$ as well.
Since $N_N(\mu_j)\le j$, we have proved the quantitative eigenvalue bound
\begin{equation}\label{eq:quantitative-eigen}
 |\Omega|\mu_j+2\norm{U_{\mu_j}-I}\le4\pi j\qquad(j\ge1).
\end{equation}
The fixed-space estimate \eqref{eq:fixed} rules out $U_E=I$ for every
$E>0$: the identity would have an infinite-dimensional fixed space.
Thus the inequalities in Theorem~\ref{thm:main} are strict for smooth
domains. The quantitative form \eqref{eq:quantitative-eigen} is needed
for the nonsmooth limit.

\section{Lipschitz domains}\label{sec:lipschitz}

For a Lipschitz Jordan domain, choose a positively oriented rescaled
arclength parametrization $\gamma:[0,L]\to\partial\Omega$, based at zero.
Equation~\eqref{eq:transport} has a unitary, absolutely continuous
solution, since its coefficient is in $L^\infty$.
Orientation-preserving bi-Lipschitz reparametrizations do not change
$V_E$. Changing the base point conjugates $V_E$ by a unitary transport
matrix. In particular $\norm{U_E-I}$ is independent of these choices.

\begin{lemma}\label{lem:lipschitz-fixed}
For every bounded simply connected Lipschitz domain and every $E>0$,
\begin{equation}\label{eq:lipschitz-fixed}
 \dim\ker(V_E-I)\le\dim\ker(-\Delta_\Omega^N-E)<\infty.
\end{equation}
\end{lemma}
\begin{proof}
The observation $\OO_E$ is injective also for this parametrization.
Indeed, if $f=W_Ev$ has $f_0=0$, the first transport row and
$|\gamma'|>0$ almost everywhere imply $f_1=0$ almost everywhere;
absolute continuity and the successive rows give $f_n=0$ for every $n$.

If $V_Ev=v$, then $f=W_Ev$ is periodic and $f_0\in W^{1,\infty}(\T)$.
The Herglotz driven identities hold almost everywhere on the Lipschitz
boundary. Pairing $f$ with a driven vector and integrating gives
$\ip{f_0}{g_a}=0$ for every Herglotz density $a$.
The outgoing double-layer potential with density $f_0$ has zero far
field, by conjugation and reversal of plane-wave directions. Rellich
uniqueness and the connectedness of the exterior make it zero outside.
The negative of this double-layer potential is an interior $H^1$
solution with zero weak Neumann data and Dirichlet trace $f_0$, by
the Lipschitz mapping and jump relations in \cite{McLean}.
The trace identifies this map injectively, by injectivity of $\OO_E$.
The compact embedding $H^1(\Omega)\hookrightarrow L^2(\Omega)$ makes
the Neumann eigenspace finite-dimensional, proving
\eqref{eq:lipschitz-fixed}. Since $\HH$ is infinite-dimensional,
this also proves $U_E\ne I$.
\end{proof}

\begin{lemma}\label{lem:curve-continuity}
Let $\gamma_n,\gamma$ be closed Lipschitz curves on $[0,L]$, all based
at zero, with $\gamma_n\to\gamma$ uniformly and
$\sup_n\int_0^L|\gamma_n'|<\infty$. If $E_n\to E>0$, then their
transport monodromies satisfy
\begin{equation}\label{eq:curve-continuity}
 \norm{U_{\gamma_n}(E_n)-U_\gamma(E)}\longrightarrow0.
\end{equation}
\end{lemma}
\begin{proof}
Write $\gamma_n=x_n+iy_n$, $k_n=\sqrt{E_n}$, and use the bounded
skew-adjoint operators $G_x,G_y$ from Lemma~\ref{lem:area}. Set
\[
 C_n=k_n(x_n'G_x+y_n'G_y),\quad C=k(x'G_x+y'G_y),\quad
 F_n=k(xG_x+yG_y)-k_n(x_nG_x+y_nG_y).
\]
Then $F_n'=C-C_n$, $F_n(0)=F_n(L)=0$, and
$\norm{F_n}_\infty\to0$. If $W_n'=C_nW_n$ and $W'=CW$, integration
by parts in the absolutely continuous operator-valued functions gives
\begin{align*}
 V_n^*V-I
 &=\int_0^L W_n^*F_n'W\,d\theta\\
 &=\int_0^L W_n^*(C_nF_n-F_nC)W\,d\theta.
\end{align*}
Unitarity yields
\[
 \norm{V_n^*V-I}\le\norm{F_n}_\infty
    \int_0^L(\norm{C_n}+\norm C)\,d\theta\longrightarrow0.
\]
This proves \eqref{eq:curve-continuity}. In particular, convergence
of the derivatives $\gamma_n'$ is not required.
\end{proof}

We now apply smooth approximation. Neumann domain monotonicity is
neither assumed nor used. Theorem~5.1 and Remark~5.3 of
Ball--Zarnescu \cite{BZ} provide smooth domains $\Omega_\epsilon$ and
uniformly bi-Lipschitz maps $f_\epsilon:\overline\Omega\to
\overline{\Omega_\epsilon}$, preserving topology and equal to the identity
outside a boundary collar shrinking to zero. In particular
$\Omega_\epsilon$ is simply connected. Here $\epsilon\downarrow0$ may
be chosen for the inner approximation. The same construction is jointly
continuous at $\epsilon=0$, so
$\sup_{\overline\Omega}|f_\epsilon-\operatorname{id}|\to0$.
Pulling the Neumann form back to
$\Omega$ gives
\[
 q_\epsilon[u]=\int_\Omega A_\epsilon\nabla u\cdot\nabla\overline u,
 \qquad m_\epsilon[u]=\int_\Omega\rho_\epsilon|u|^2,
\]
where
$\rho_\epsilon=|\det Df_\epsilon|$ and
$A_\epsilon=\rho_\epsilon Df_\epsilon^{-1}Df_\epsilon^{-T}$.
They are uniformly bounded and uniformly positive; they equal $1$ and
$I$, respectively, off the shrinking collar.
For each fixed finite-dimensional subspace of $H^1(\Omega)$, dominated
convergence of these forms and min--max give
$\limsup\mu_j(\Omega_\epsilon)\le\mu_j(\Omega)$.
Conversely, pull back the first $j+1$ orthonormal eigenfunctions.
They are bounded in $H^1(\Omega)$ and have a subsequence converging
weakly in $H^1$ and strongly in $L^2$. Weighted orthonormality passes
to ordinary orthonormality, since $\rho_\epsilon$ is bounded and tends
to one almost everywhere. On each fixed interior subdomain the energy
is eventually the ordinary Dirichlet energy, so weak lower
semicontinuity followed by exhaustion gives the required lower energy
bound for every linear combination of the limits. More explicitly,
after taking a subsequence realizing the lower limit, the limiting
functions $u_0,\ldots,u_j$ are orthonormal and satisfy
\[
 q\!\left[\sum_{\ell=0}^j a_\ell u_\ell\right]
 \le \left(\liminf_{\epsilon\downarrow0}
                  \mu_j(\Omega_\epsilon)\right)
                  \sum_{\ell=0}^j|a_\ell|^2.
\]
The min--max principle gives
$\mu_j(\Omega)\le\liminf\mu_j(\Omega_\epsilon)$.
Thus
\[
 \mu_j(\Omega_\epsilon)\longrightarrow\mu_j(\Omega),\qquad
 |\Omega_\epsilon|=\int_\Omega\rho_\epsilon\longrightarrow|\Omega|.
\]
For the boundary transport use
\[
 \gamma_\epsilon=f_\epsilon\circ\gamma-f_\epsilon(\gamma(0)).
\]
These curves parametrize translated smooth boundaries with positive
orientation. Their lengths are uniformly bounded, by the uniform
Lipschitz bound on $f_\epsilon$, and they converge uniformly to $\gamma$.
Although these particular parametrizations need not be smooth,
reparametrization invariance identifies their monodromies with the
ones used in the smooth-domain proof. Lemma~\ref{lem:curve-continuity}
and eigenvalue convergence imply
\[
 U_{\gamma_\epsilon}(\mu_j(\Omega_\epsilon))
       \longrightarrow U_\gamma(\mu_j(\Omega))
       \quad\text{in operator norm}.
\]
Passing \eqref{eq:quantitative-eigen}, rather than only its strict
consequence, to the limit gives
\begin{equation}\label{eq:lipschitz-remainder}
 |\Omega|\mu_j(\Omega)
    +2\norm{U_\gamma(\mu_j(\Omega))-I}\le4\pi j\qquad(j\ge1).
\end{equation}
The additional term in \eqref{eq:lipschitz-remainder} is positive
by Lemma~\ref{lem:lipschitz-fixed}.
Thus $|\Omega|\mu_j(\Omega)<4\pi j$. If $m=N_N(E)$, then
$E\le\mu_m(\Omega)<4\pi m/|\Omega|$, proving the strict counting
inequality as well. This completes the proof of Theorem~\ref{thm:main}.

\begin{thebibliography}{99}

\bibitem{BZ}
J.~M. Ball and A. Zarnescu,
\emph{Partial regularity and smooth topology-preserving approximations
of rough domains}, Calc. Var. Partial Differential Equations
\textbf{56} (2017), article 13.

\bibitem{CK}
D. Colton and R. Kress,
\emph{On the denseness of Herglotz wave functions and electromagnetic
Herglotz pairs in Sobolev spaces}, Math. Methods Appl. Sci.
\textbf{24} (2001), 1289--1303.

\bibitem{EP}
J.-P. Eckmann and C.-A. Pillet,
\emph{Zeta functions with Dirichlet and Neumann boundary conditions
for exterior domains}, Helv. Phys. Acta \textbf{70} (1997), 44--65.

\bibitem{Filonov}
N. Filonov,
\emph{On the P\'olya conjecture for the Neumann problem in planar convex
domains}, Comm. Pure Appl. Math. \textbf{78} (2025), 537--544.

\bibitem{FLPS}
N. Filonov, M. Levitin, I. Polterovich, and D.~A. Sher,
\emph{P\'olya's conjecture for Euclidean balls}, Invent. Math.
\textbf{234} (2023), 129--169.

\bibitem{FLPS26}
N. Filonov, M. Levitin, I. Polterovich, and D.~A. Sher,
\emph{P\'olya's conjecture for higher-dimensional Neumann balls},
arXiv:2607.29305, 2026.

\bibitem{FrankLarson}
R.~L. Frank and S. Larson,
\emph{Semiclassical inequalities for Dirichlet and Neumann Laplacians
on convex domains}, Comm. Pure Appl. Math. \textbf{79} (2026), 762--822.

\bibitem{FLP}
P. Freitas, J. Lagac\'e, and J. Payette,
\emph{Optimal unions of scaled copies of domains and P\'olya's conjecture},
Ark. Mat. \textbf{59} (2021), 11--51.

\bibitem{FS}
P. Freitas and I. Salavessa,
\emph{Families of non-tiling domains satisfying P\'olya's conjecture},
J. Math. Phys. \textbf{64} (2023), article 121503.

\bibitem{HW}
X. He and Z. Wang,
\emph{P\'olya's conjecture for thin products},
Int. Math. Res. Not. \textbf{2026} (2026), no.~16, article rnag182.

\bibitem{Kellner}
R. Kellner, \emph{On a theorem of Polya}, Amer. Math. Monthly
\textbf{73} (1966), 856--858.

\bibitem{Kroger}
P. Kr\"oger,
\emph{Upper bounds for the Neumann eigenvalues on a bounded domain
in Euclidean space}, J. Funct. Anal. \textbf{106} (1992), 353--357.

\bibitem{Laptev}
A. Laptev,
\emph{Dirichlet and Neumann eigenvalue problems on domains in Euclidean
spaces}, J. Funct. Anal. \textbf{151} (1997), 531--545.

\bibitem{Li26}
Y. Li,
\emph{P\'olya's conjecture for the Neumann eigenvalues on Euclidean balls},
arXiv:2607.25958, 2026.

\bibitem{LM}
J.-L. Lions and E. Magenes,
\emph{Non-Homogeneous Boundary Value Problems and Applications},
Vol.~I, Grundlehren der mathematischen Wissenschaften, vol.~181,
Springer-Verlag, Berlin, 1972.

\bibitem{McLean}
W. McLean, \emph{Strongly Elliptic Systems and Boundary Integral Equations},
Cambridge University Press, Cambridge, 2000.

\bibitem{MHP}
A. Moiola, R. Hiptmair, and I. Perugia,
\emph{Vekua theory for the Helmholtz operator}, Z. Angew. Math. Phys.
\textbf{62} (2011), 779--807.

\bibitem{Polya54}
G. P\'olya, \emph{Mathematics and Plausible Reasoning},
Vols.~I--II, Princeton University Press, Princeton, NJ, 1954.

\bibitem{Polya}
G. P\'olya, \emph{On the eigenvalues of vibrating membranes},
Proc. London Math. Soc. (3) \textbf{11} (1961), 419--433.

\bibitem{Pommerenke}
Ch. Pommerenke, \emph{Boundary Behaviour of Conformal Maps},
Grundlehren der mathematischen Wissenschaften, vol.~299,
Springer-Verlag, Berlin, 1992.

\bibitem{Rudin}
W. Rudin, \emph{Real and Complex Analysis}, 3rd ed.,
McGraw--Hill, New York, 1987.

\bibitem{Simon}
B. Simon, \emph{Trace Ideals and Their Applications}, 2nd ed.,
Mathematical Surveys and Monographs, vol.~120,
American Mathematical Society, Providence, RI, 2005.

\bibitem{Szego}
G. Szeg\H{o},
\emph{Inequalities for certain eigenvalues of a membrane of given area},
J. Rational Mech. Anal. \textbf{3} (1954), 343--356.

\bibitem{Vekua}
I.~N. Vekua, \emph{New Methods for Solving Elliptic Equations},
North-Holland Publishing Co., Amsterdam, 1967.

\bibitem{Weinberger}
H.~F. Weinberger,
\emph{An isoperimetric inequality for the $N$-dimensional free membrane
problem}, J. Rational Mech. Anal. \textbf{5} (1956), 633--636.

\bibitem{Weyl}
H. Weyl, \emph{\"Uber die asymptotische Verteilung der Eigenwerte},
Nachr. Ges. Wiss. G\"ottingen, Math.-Phys. Kl. (1911), 110--117.

\end{thebibliography}
\end{document}
